% GENERATED FILE. Edit canonical lessons, then run npm run generate:books.
% canonical-source-hash: fnv1a64:858db286c418ed21
% canonical-lessons: LA-C496-paveo, LA-C496-vereor, LA-C496-perterreo, LA-C496-horribilis, LA-C496-terribilis, LA-R496-review-489-496

\chapter{Fear and dread}
\label{ch:la-fear-and-dread}
\hlchaptermodality{496}

\begin{chapteropening}
\textbf{By the end of this chapter:} \emph{I can describe fear and dreadful things.}

Complete the last lesson of chapter 496: I can describe fear and dreadful things.
\end{chapteropening}

\section[paveō, pavēre]{paveō, pavēre — to quake with fear, be terrified}

\label{lesson:LA-C496-paveo}



\begin{quote}
Before the new one: say the Latin for thunderstruck, then the Latin for to be struck dumb.
\end{quote}

\subsection*{You'll want to know: paveō, pavēre}
\textbf{paveō, pavēre} — \textquotedblleft{}to quake with fear, be terrified\textquotedblright{}.

\begin{quote}
Ovēs lupum vident et pavent.
\end{quote}

\begin{quote}
\emph{The sheep catch sight of the wolf and quake with fear.}
\end{quote}

\emph{paveō, pavēre, pāvī}. The noun is \emph{pavor} — terror.

\begin{cousinweb}
\textbf{paveō} is a native Latin verb; EtymDB links it with a Proto-Indo-European root meaning \textquotedblleft{}to strike\textquotedblright{}, but Lewis \& Short calls its origin doubtful, so it is not settled.
\end{cousinweb}

\begin{grammarlens}[title={What you've built}]
One more word for telling, explaining and describing at B1.
\end{grammarlens}

\subsection*{Your turn}
\begin{itemize}
\raggedright
  \item \emph{Say it:} \emph{paveō, pavēre}
  \item \emph{Say it:} \emph{paveō, pavēre}, once more
  \item \emph{Say it:} a sentence of your own with \emph{paveō, pavēre}
  \item \emph{Recall:} say the Latin for thunderstruck, then the Latin for to be struck dumb, then say \emph{paveō, pavēre} again
\end{itemize}

\begin{tcolorbox}[breakable,colback=teal!4,colframe=teal!35!black,title={Before you move on}]
What does paveō, pavēre mean? (\textquotedblleft{}To quake with fear, be terrified\textquotedblright{}.) Say it once more.
\end{tcolorbox}
\section[vereor, verērī]{vereor, verērī — to respect; to fear}

\label{lesson:LA-C496-vereor}



\begin{quote}
Before the new one: say the Latin for to be struck dumb, then the Latin for to quake with fear.
\end{quote}

\subsection*{You'll want to know: vereor, verērī}
\textbf{vereor, verērī} — \textquotedblleft{}to respect; to fear\textquotedblright{}.

\begin{quote}
Discipulī magistrum verentur.
\end{quote}

\begin{quote}
\emph{The pupils respect their teacher.}
\end{quote}

Deponent: \emph{vereor, verērī, veritus sum}. Fear mixed with respect. \emph{Vereor nē…} — I am afraid that…

\begin{cousinweb}
\textbf{vereor} is inherited through Proto-Italic from a Proto-Indo-European root \emph{wer-} (EtymDB), and Lewis \& Short compares Greek \emph{horao} \textquotedblleft{}to see\textquotedblright{} and English \emph{ward}. English \emph{revere} comes from \emph{revereor}.
\end{cousinweb}

\begin{grammarlens}[title={What you've built}]
One more word for telling, explaining and describing at B1.
\end{grammarlens}

\subsection*{Your turn}
\begin{itemize}
\raggedright
  \item \emph{Say it:} \emph{vereor, verērī}
  \item \emph{Say it:} \emph{vereor, verērī}, once more
  \item \emph{Say it:} a sentence of your own with \emph{vereor, verērī}
  \item \emph{Recall:} say the Latin for to be struck dumb, then the Latin for to quake with fear, then say \emph{vereor, verērī} again
\end{itemize}

\begin{tcolorbox}[breakable,colback=teal!4,colframe=teal!35!black,title={Before you move on}]
What does vereor, verērī mean? (\textquotedblleft{}To respect; to fear\textquotedblright{}.) Say it once more.
\end{tcolorbox}
\section[perterreō, perterrēre]{perterreō, perterrēre — to terrify thoroughly}

\label{lesson:LA-C496-perterreo}



\begin{quote}
Before the new one: say the Latin for to quake with fear, then the Latin for to respect.
\end{quote}

\subsection*{You'll want to know: perterreō, perterrēre}
\textbf{perterreō, perterrēre} — \textquotedblleft{}to terrify thoroughly\textquotedblright{}.

\begin{quote}
Tonitrus puerōs perterret.
\end{quote}

\begin{quote}
\emph{The thunder terrifies the children.}
\end{quote}

\emph{perterreō, perterrēre, perterruī, perterritum}: \emph{per-} 'thoroughly' + \emph{terreō} 'to frighten'.

\begin{cousinweb}
\textbf{perterreō} is built on \emph{per-} \textquotedblleft{}thoroughly\textquotedblright{} and \emph{terreō} \textquotedblleft{}to frighten\textquotedblright{} (Lewis \& Short; EtymDB). English \emph{terrify} and \emph{terror} belong to the same family.
\end{cousinweb}

\begin{grammarlens}[title={What you've built}]
One more word for telling, explaining and describing at B1.
\end{grammarlens}

\subsection*{Your turn}
\begin{itemize}
\raggedright
  \item \emph{Say it:} \emph{perterreō, perterrēre}
  \item \emph{Say it:} \emph{perterreō, perterrēre}, once more
  \item \emph{Say it:} a sentence of your own with \emph{perterreō, perterrēre}
  \item \emph{Recall:} say the Latin for to quake with fear, then the Latin for to respect, then say \emph{perterreō, perterrēre} again
\end{itemize}

\begin{tcolorbox}[breakable,colback=teal!4,colframe=teal!35!black,title={Before you move on}]
What does perterreō, perterrēre mean? (\textquotedblleft{}To terrify thoroughly\textquotedblright{}.) Say it once more.
\end{tcolorbox}
\section[horribilis, horribile]{horribilis, horribile — dreadful, horrible}

\label{lesson:LA-C496-horribilis}



\begin{quote}
Before the new one: say the Latin for to respect, then the Latin for to terrify thoroughly.
\end{quote}

\subsection*{You'll want to know: horribilis, horribile}
\textbf{horribilis, horribile} — \textquotedblleft{}dreadful, horrible\textquotedblright{}.

\begin{quote}
Tempestās horribilis nāvem frēgit.
\end{quote}

\begin{quote}
\emph{A dreadful storm broke the ship.}
\end{quote}

Two endings. From \emph{horreō} — to bristle, to shudder. English \emph{horrible}.

\begin{cousinweb}
\textbf{horribilis} is built on \emph{horreō} \textquotedblleft{}to bristle, to shudder\textquotedblright{} with \emph{-ibilis} (Lewis \& Short; EtymDB). English \emph{horrible} comes from it.
\end{cousinweb}

\begin{grammarlens}[title={What you've built}]
One more word for telling, explaining and describing at B1.
\end{grammarlens}

\subsection*{Your turn}
\begin{itemize}
\raggedright
  \item \emph{Say it:} \emph{horribilis, horribile}
  \item \emph{Say it:} \emph{horribilis, horribile}, once more
  \item \emph{Say it:} a sentence of your own with \emph{horribilis, horribile}
  \item \emph{Recall:} say the Latin for to respect, then the Latin for to terrify thoroughly, then say \emph{horribilis, horribile} again
\end{itemize}

\begin{tcolorbox}[breakable,colback=teal!4,colframe=teal!35!black,title={Before you move on}]
What does horribilis, horribile mean? (\textquotedblleft{}Dreadful, horrible\textquotedblright{}.) Say it once more.
\end{tcolorbox}
\section[terribilis, terribile]{terribilis, terribile — frightful, terrible}

\label{lesson:LA-C496-terribilis}



\begin{quote}
Before the new one: say the Latin for to terrify thoroughly, then the Latin for dreadful.
\end{quote}

\subsection*{You'll want to know: terribilis, terribile}
\textbf{terribilis, terribile} — \textquotedblleft{}frightful, terrible\textquotedblright{}.

\begin{quote}
Leō vōcem terribilem habet.
\end{quote}

\begin{quote}
\emph{The lion has a terrible voice.}
\end{quote}

Two endings. From \emph{terreō} — to frighten. English \emph{terrible}.

\begin{cousinweb}
\textbf{terribilis} is built on \emph{terreō} \textquotedblleft{}to frighten\textquotedblright{} with \emph{-ibilis} (Lewis \& Short); EtymDB traces \emph{terreō} to a Proto-Indo-European root meaning \textquotedblleft{}to tremble, to frighten\textquotedblright{}. English \emph{terrible} comes from it.
\end{cousinweb}

\begin{grammarlens}[title={What you've built}]
One more word for telling, explaining and describing at B1.
\end{grammarlens}

\subsection*{Your turn}
\begin{itemize}
\raggedright
  \item \emph{Say it:} \emph{terribilis, terribile}
  \item \emph{Say it:} \emph{terribilis, terribile}, once more
  \item \emph{Say it:} a sentence of your own with \emph{terribilis, terribile}
  \item \emph{Recall:} say the Latin for to terrify thoroughly, then the Latin for dreadful, then say \emph{terribilis, terribile} again
\end{itemize}

\begin{tcolorbox}[breakable,colback=teal!4,colframe=teal!35!black,title={Before you move on}]
What does terribilis, terribile mean? (\textquotedblleft{}Frightful, terrible\textquotedblright{}.) Say it once more.
\end{tcolorbox}
\section[(dialogue)]{Feelings and character}

\label{lesson:LA-R496-review-489-496}



\begin{quote}
Say the Latin for joy and to rejoice.
\end{quote}

\subsection*{Your turn}
Say these aloud:

\begin{itemize}
\raggedright
  \item joy, to rejoice, cheerfulness, to leap for joy, pleasantness
  \item sadness, sorrow, to grieve, grief, sorrowful
  \item to weep, to wail, to groan, a groan, to sigh
  \item to console, consolation, solace, gentle, to soothe
  \item to grow angry, hot-tempered, a hot temper, unworthy, to be indignant
  \item to hate, hateful, to turn up one's nose at, disgust, it wearies
  \item admiration, to wonder, to be stunned, thunderstruck, to be struck dumb
  \item to quake with fear, to respect, to terrify thoroughly, dreadful, frightful
\end{itemize}

\begin{tcolorbox}[breakable,colback=teal!4,colframe=teal!35!black,title={Before you move on}]
Joy? (\textbf{laetitia, laetitiae}.) To rejoice? (\textbf{laetor, laetārī}.) Cheerfulness? (\textbf{hilaritās, hilaritātis}.) To leap for joy? (\textbf{exsultō, exsultāre}.) Pleasantness? (\textbf{iūcunditās, iūcunditātis}.) Sadness? (\textbf{trīstitia, trīstitiae}.) Sorrow? (\textbf{maestitia, maestitiae}.) To grieve? (\textbf{maereō, maerēre}.) Grief? (\textbf{maeror, maerōris}.) Sorrowful? (\textbf{lūctuōsus, lūctuōsa, lūctuōsum}.) To weep? (\textbf{lacrimō, lacrimāre}.) To wail? (\textbf{plōrō, plōrāre}.) To groan? (\textbf{gemō, gemere}.) A groan? (\textbf{gemitus, gemitūs}.) To sigh? (\textbf{suspīrō, suspīrāre}.) To console? (\textbf{cōnsōlor, cōnsōlārī}.) Consolation? (\textbf{cōnsōlātiō, cōnsōlātiōnis}.) Solace? (\textbf{sōlācium, sōlāciī}.) Gentle? (\textbf{lēnis, lēne}.) To soothe? (\textbf{lēniō, lēnīre}.) To grow angry? (\textbf{īrāscor, īrāscī}.) Hot-tempered? (\textbf{īrācundus, īrācunda, īrācundum}.) A hot temper? (\textbf{īrācundia, īrācundiae}.) Unworthy? (\textbf{indignus, indigna, indignum}.) To be indignant? (\textbf{indignor, indignārī}.) To hate? (\textbf{ōdī, ōdisse}.) Hateful? (\textbf{odiōsus, odiōsa, odiōsum}.) To turn up one's nose at? (\textbf{fastīdiō, fastīdīre}.) Disgust? (\textbf{fastīdium, fastīdiī}.) It wearies? (\textbf{taedet, taedēre}.) Admiration? (\textbf{admīrātiō, admīrātiōnis}.) To wonder? (\textbf{mīror, mīrārī}.) To be stunned? (\textbf{stupeō, stupēre}.) Thunderstruck? (\textbf{attonitus, attonita, attonitum}.) To be struck dumb? (\textbf{obstupēscō, obstupēscere}.) To quake with fear? (\textbf{paveō, pavēre}.) To respect? (\textbf{vereor, verērī}.) To terrify thoroughly? (\textbf{perterreō, perterrēre}.) Dreadful? (\textbf{horribilis, horribile}.) Frightful? (\textbf{terribilis, terribile}.)
\end{tcolorbox}
